% GENERATED FILE. Edit canonical lessons, then run npm run generate:books.
% canonical-source-hash: fnv1a64:d7c0642b350c1fb2
% canonical-lessons: ML-C355-patrapravarttakan, ML-C355-aropanam, ML-C355-velippetuttuka, ML-C355-ottukutuka, ML-C355-kuvuka

\chapter{Journalists and Crowds}
\label{ch:ml-journalists-and-crowds}
\hlchaptermodality{355}

\begin{chapteropening}
\textbf{By the end of this chapter:} \emph{I can say a journalist, an allegation, to reveal, to gather and to boo.}

Complete the last lesson of chapter 355: I can say a journalist, an allegation, to reveal, to gather and to boo.
\end{chapteropening}

\section[\texorpdfstring{\ml{പത്രപ്രവർത്തകൻ} (patrapravarttakan)}{പത്രപ്രവർത്തകൻ (patrapravarttakan)}]{\ml{പത്രപ്രവർത്തകൻ} (patrapravarttakan) — a journalist}

\label{lesson:ML-C355-patrapravarttakan}



\begin{quote}
Before the new one: say the Malayalam for to applaud, then the Malayalam for to mention.
\end{quote}

\subsection*{You'll want to know: \ml{പത്രപ്രവർത്തകൻ}}
\textbf{\ml{പത്രപ്രവർത്തകൻ}} — \emph{patrapravarttakan} — \textquotedblleft{}a journalist\textquotedblright{}.

\begin{quote}
\ml{എന്റെ} \ml{അമ്മാവൻ} \ml{ഒരു} \ml{പത്രപ്രവർത്തകൻ} \ml{ആണ്}.
\end{quote}

\begin{quote}
\emph{enṟe ammāvan oru patrapravarttakan āṇŭ.}
\end{quote}

\begin{quote}
My uncle is a journalist.
\end{quote}

A woman journalist is \ml{പത്രപ്രവർത്തക}. The trade is \ml{പത്രപ്രവർത്തനം}, journalism.

\begin{cousinweb}
\textbf{\ml{പത്രപ്രവർത്തകൻ}} (\emph{patrapravarttakan}) is built from \ml{പത്രം} (\emph{patraṁ}), \textquotedblleft{}a newspaper\textquotedblright{}, and \ml{പ്രവർത്തകൻ}, \textquotedblleft{}a worker\textquotedblright{}, from Sanskrit \emph{pravartaka}, \textquotedblleft{}setting in motion, promoting\textquotedblright{} (Monier-Williams), the same word as in \ml{സഹപ്രവർത്തകൻ}, \textquotedblleft{}a colleague\textquotedblright{}.
\end{cousinweb}

\begin{grammarlens}[title={What you've built}]
One more word for telling what happened.
\end{grammarlens}

\subsection*{Your turn}
\begin{itemize}
\raggedright
  \item \emph{Say it:} \emph{patrapravarttakan}
  \item \emph{Say it:} \emph{patrapravarttakan}, once more
  \item \emph{Say it:} \emph{enṟe ammāvan oru patrapravarttakan āṇŭ.} — My uncle is a journalist.
  \item \emph{Recall:} say the Malayalam for to applaud, then the Malayalam for to mention, then say \emph{patrapravarttakan} again
\end{itemize}

\begin{tcolorbox}[breakable,colback=teal!4,colframe=teal!35!black,title={Before you move on}]
What does \ml{പത്രപ്രവർത്തകൻ} mean? (\textquotedblleft{}A journalist\textquotedblright{}.) Say it once more.
\end{tcolorbox}
\section[\texorpdfstring{\ml{ആരോപണം} (ārōpaṇaṁ)}{ആരോപണം (ārōpaṇaṁ)}]{\ml{ആരോപണം} (ārōpaṇaṁ) — an allegation; a charge}

\label{lesson:ML-C355-aropanam}



\begin{quote}
Before the new one: say the Malayalam for to mention, then the Malayalam for a journalist.
\end{quote}

\subsection*{You'll want to know: \ml{ആരോപണം}}
\textbf{\ml{ആരോപണം}} — \emph{ārōpaṇaṁ} — \textquotedblleft{}an allegation; a charge\textquotedblright{}.

\begin{quote}
\ml{അവന്} \ml{എതിരെ} \ml{ഒരു} \ml{ആരോപണം} \ml{ഉണ്ട്}.
\end{quote}

\begin{quote}
\emph{avanŭ etire oru ārōpaṇaṁ uṇṭŭ.}
\end{quote}

\begin{quote}
There is an allegation against him.
\end{quote}

Against takes \ml{എതിരെ}: \ml{അവന്} \ml{എതിരെ}, against him.

\begin{cousinweb}
\textbf{\ml{ആരോപണം}} (\emph{ārōpaṇaṁ}) is borrowed from Sanskrit \emph{āropaṇa}, \textquotedblleft{}causing to mount, raising up\textquotedblright{} (Monier-Williams). An allegation lays a charge upon someone.
\end{cousinweb}

\begin{grammarlens}[title={What you've built}]
One more word for telling what happened.
\end{grammarlens}

\subsection*{Your turn}
\begin{itemize}
\raggedright
  \item \emph{Say it:} \emph{ārōpaṇaṁ}
  \item \emph{Say it:} \emph{ārōpaṇaṁ}, once more
  \item \emph{Say it:} \emph{avanŭ etire oru ārōpaṇaṁ uṇṭŭ.} — There is an allegation against him.
  \item \emph{Recall:} say the Malayalam for to mention, then the Malayalam for a journalist, then say \emph{ārōpaṇaṁ} again
\end{itemize}

\begin{tcolorbox}[breakable,colback=teal!4,colframe=teal!35!black,title={Before you move on}]
What does \ml{ആരോപണം} mean? (\textquotedblleft{}An allegation; a charge\textquotedblright{}.) Say it once more.
\end{tcolorbox}
\section[\texorpdfstring{\ml{വെളിപ്പെടുത്തുക} (veḷippeṭuttuka)}{വെളിപ്പെടുത്തുക (veḷippeṭuttuka)}]{\ml{വെളിപ്പെടുത്തുക} (veḷippeṭuttuka) — to reveal; to disclose}

\label{lesson:ML-C355-velippetuttuka}



\begin{quote}
Before the new one: say the Malayalam for a journalist, then the Malayalam for an allegation.
\end{quote}

\subsection*{You'll want to know: \ml{വെളിപ്പെടുത്തുക}}
\textbf{\ml{വെളിപ്പെടുത്തുക}} — \emph{veḷippeṭuttuka} — \textquotedblleft{}to reveal; to disclose\textquotedblright{}.

\begin{quote}
\ml{ലേഖകൻ} \ml{സംഭവത്തിന്റെ} \ml{സത്യാവസ്ഥ} \ml{വെളിപ്പെടുത്തി}.
\end{quote}

\begin{quote}
\emph{lēkhakan sambhavattinṟe satyāvastha veḷippeṭutti.}
\end{quote}

\begin{quote}
The correspondent revealed the real facts of the incident.
\end{quote}

Past: \ml{വെളിപ്പെടുത്തി}, \emph{veḷippeṭutti}.

Built on \ml{വെളി}, the open, the outside: to bring into the open.

\begin{cousinweb}
\textbf{\ml{വെളിപ്പെടുത്തുക}} (\emph{veḷippeṭuttuka}) is native Dravidian: \ml{വെളി} (\emph{veḷi}), \textquotedblleft{}an open field; the outside\textquotedblright{} (DEDR 5498), with -\ml{പ്പെടുത്തുക}, \textquotedblleft{}to cause to be\textquotedblright{}. To reveal is to put a thing out in the open.
\end{cousinweb}

\begin{grammarlens}[title={What you've built}]
One more word for telling what happened.
\end{grammarlens}

\subsection*{Your turn}
\begin{itemize}
\raggedright
  \item \emph{Say it:} \emph{veḷippeṭuttuka}
  \item \emph{Say it:} \emph{veḷippeṭuttuka}, once more
  \item \emph{Say it:} \emph{lēkhakan sambhavattinṟe satyāvastha veḷippeṭutti.} — The correspondent revealed the real facts of the incident.
  \item \emph{Recall:} say the Malayalam for a journalist, then the Malayalam for an allegation, then say \emph{veḷippeṭuttuka} again
\end{itemize}

\begin{tcolorbox}[breakable,colback=teal!4,colframe=teal!35!black,title={Before you move on}]
What does \ml{വെളിപ്പെടുത്തുക} mean? (\textquotedblleft{}To reveal; to disclose\textquotedblright{}.) Say it once more.
\end{tcolorbox}
\section[\texorpdfstring{\ml{ഒത്തുകൂടുക} (ottukūṭuka)}{ഒത്തുകൂടുക (ottukūṭuka)}]{\ml{ഒത്തുകൂടുക} (ottukūṭuka) — to gather; to assemble}

\label{lesson:ML-C355-ottukutuka}



\begin{quote}
Before the new one: say the Malayalam for an allegation, then the Malayalam for to reveal.
\end{quote}

\subsection*{You'll want to know: \ml{ഒത്തുകൂടുക}}
\textbf{\ml{ഒത്തുകൂടുക}} — \emph{ottukūṭuka} — \textquotedblleft{}to gather; to assemble\textquotedblright{}.

\begin{quote}
\ml{വാർത്ത} \ml{കേട്ട്} \ml{ആളുകൾ} \ml{കവലയിൽ} \ml{ഒത്തുകൂടി}.
\end{quote}

\begin{quote}
\emph{vārtta kēṭṭŭ āḷukaḷ kavalayil ottukūṭi.}
\end{quote}

\begin{quote}
Hearing the news, people gathered at the junction.
\end{quote}

Past: \ml{ഒത്തുകൂടി}, \emph{ottukūṭi}.

The subject is plural: people gather. \ml{ഒത്ത്} means together.

\begin{cousinweb}
\textbf{\ml{ഒത്തുകൂടുക}} (\emph{ottukūṭuka}) is native Dravidian: \ml{കൂടുക} (\emph{kūṭuka}), \textquotedblleft{}to come together, to meet\textquotedblright{} (DEDR 1882), with \ml{ഒത്ത്}, \textquotedblleft{}together\textquotedblright{}, in front.
\end{cousinweb}

\begin{grammarlens}[title={What you've built}]
One more word for telling what happened.
\end{grammarlens}

\subsection*{Your turn}
\begin{itemize}
\raggedright
  \item \emph{Say it:} \emph{ottukūṭuka}
  \item \emph{Say it:} \emph{ottukūṭuka}, once more
  \item \emph{Say it:} \emph{vārtta kēṭṭŭ āḷukaḷ kavalayil ottukūṭi.} — Hearing the news, people gathered at the junction.
  \item \emph{Recall:} say the Malayalam for an allegation, then the Malayalam for to reveal, then say \emph{ottukūṭuka} again
\end{itemize}

\begin{tcolorbox}[breakable,colback=teal!4,colframe=teal!35!black,title={Before you move on}]
What does \ml{ഒത്തുകൂടുക} mean? (\textquotedblleft{}To gather; to assemble\textquotedblright{}.) Say it once more.
\end{tcolorbox}
\section[\texorpdfstring{\ml{കൂവുക} (kūvuka)}{കൂവുക (kūvuka)}]{\ml{കൂവുക} (kūvuka) — to boo; to hoot}

\label{lesson:ML-C355-kuvuka}



\begin{quote}
Before the new one: say the Malayalam for to reveal, then the Malayalam for to gather.
\end{quote}

\subsection*{You'll want to know: \ml{കൂവുക}}
\textbf{\ml{കൂവുക}} — \emph{kūvuka} — \textquotedblleft{}to boo; to hoot\textquotedblright{}.

\begin{quote}
\ml{കളി} \ml{തോറ്റപ്പോൾ} \ml{ആളുകൾ} \ml{കൂവി}.
\end{quote}

\begin{quote}
\emph{kaḷi tōṟṟappōḷ āḷukaḷ kūvi.}
\end{quote}

\begin{quote}
When the game was lost, people booed.
\end{quote}

Past: \ml{കൂവി}, \emph{kūvi}.

The same verb is the crowing of a cock at dawn: \ml{കോഴി} \ml{കൂവി}. At a speech it is a jeer.

\begin{cousinweb}
\textbf{\ml{കൂവുക}} (\emph{kūvuka}) is native Dravidian: DEDR 1868 groups it with Tamil \emph{kū}, a cry, a call. Malayalam uses it for a cock crowing and for a crowd jeering.
\end{cousinweb}

\begin{grammarlens}[title={What you've built}]
One more word for telling what happened.
\end{grammarlens}

\subsection*{Your turn}
\begin{itemize}
\raggedright
  \item \emph{Say it:} \emph{kūvuka}
  \item \emph{Say it:} \emph{kūvuka}, once more
  \item \emph{Say it:} \emph{kaḷi tōṟṟappōḷ āḷukaḷ kūvi.} — When the game was lost, people booed.
  \item \emph{Recall:} say the Malayalam for to reveal, then the Malayalam for to gather, then say \emph{kūvuka} again
\end{itemize}

\begin{tcolorbox}[breakable,colback=teal!4,colframe=teal!35!black,title={Before you move on}]
What does \ml{കൂവുക} mean? (\textquotedblleft{}To boo; to hoot\textquotedblright{}.) Say it once more.
\end{tcolorbox}
